%=====================================================================
\section{The RL algorithm (PPO)}\label{sec:ppo}
%=====================================================================

We use \emph{Proximal Policy Optimization} (PPO) \cite{schulman2017ppo},
one of the most widely used RL algorithms. It is an actor--critic method
built directly on the policy gradient of Eq.~\eqref{eq:pg}. It is
implemented from scratch in about 150 lines of PyTorch
(\code{train\_ppo.py}), so every step below can be traced in the code.

\subsection{The idea}
PPO repeats one cycle:
\begin{enumerate}\itemsep1pt
\item \textbf{Collect experience.} Run the current policy (with its
exploration noise) in the simulator for a while and record every
$(o_t,a_t,r_t)$.
\item \textbf{Score each decision.} Use the critic to estimate each
action's advantage: was it better or worse than expected?
\item \textbf{Improve the policy, carefully.} Change $\theta$ to make
good actions more likely and bad ones less likely, \emph{but not by too
much in one go}.
\item \textbf{Improve the critic.} Fit $V_\phi$ to the returns just
observed, so the next round's advantage estimates are better.
\end{enumerate}
The phrase ``not by too much'' is what distinguishes PPO. A policy-gradient
step uses data collected with the \emph{old} policy, so it is only
trustworthy for policies close to it. A large step can jump to a policy
that behaves very differently, where the old data say nothing, and
performance can collapse. Engineers will recognize the same concern behind
step-size control in Newton's method or the damping term in
Levenberg--Marquardt: trust the local model only within a limited region.

\subsection{The networks}
Two separate small networks are used. Each has two hidden layers of 64
units with $\tanh$ activations.
\begin{itemize}\itemsep1pt
\item \textbf{Actor:} $o_t\mapsto\mu_\theta(o_t)$, the mean of the action.
The action is drawn as $a_t\sim\mathcal N(\mu_\theta(o_t),\sigma^2)$. The
spread $\sigma$ is one more learned parameter, initialized at
$e^{-0.5}\approx0.6$ (on an action range of $[-1,1]$, so exploration starts
wide). The environment clips the action to $[-1,1]$ and maps it to a
release.
\item \textbf{Critic:} $o_t\mapsto V_\phi(o_t)$, the predicted
(discounted) return still to come.
\end{itemize}
For evaluation, the trained actor is used deterministically:
$a_t=\mathrm{clip}(\mu_\theta(o_t))$.

\subsection{Step 2 in detail: generalized advantage estimation}
The TD error of Eq.~\eqref{eq:td} is a one-day look-ahead. It has low
variance, but it is biased when the critic is wrong. The full remaining
return is unbiased but noisy. \emph{Generalized advantage estimation} (GAE)
\cite{schulman2016gae} blends all look-ahead lengths with a geometric
weight $\lambda\in[0,1]$:
\begin{equation}
\hat A_t=\sum_{l\ge0}(\gamma\lambda)^l\,\delta_{t+l},\qquad
\delta_t=r_t+\gamma V_\phi(o_{t+1})(1-d_t)-V_\phi(o_t),
\label{eq:gae}
\end{equation}
where $d_t=1$ on the last day of a season (there is no ``tomorrow'' to
bootstrap from). $\lambda=0$ gives the one-day TD error. $\lambda=1$ gives
the full remaining return minus the critic's prediction. We use
$\lambda=0.95$, so credit for a flood is spread back over roughly the
previous $1/(1-\gamma\lambda)\approx17$ days. The sum is computed in one
backward pass over the recorded data. The critic's regression target is
$\hat R_t=\hat A_t+V_\phi(o_t)$. Before each update the advantages are
standardized (zero mean, unit variance), which only rescales the step.

\subsection{Step 3 in detail: the clipped objective}
After the data are collected, the policy used to collect them is frozen as
$\pi_{\theta_{old}}$. For a candidate $\theta$, the \emph{probability
ratio}
\[
\rho_t(\theta)=\frac{\pi_\theta(a_t\mid o_t)}{\pi_{\theta_{old}}(a_t\mid o_t)}
\]
says how much more (or less) likely the new policy is to take the recorded
action. The objective $\mathbb E[\rho_t\hat A_t]$ has the same gradient as
Eq.~\eqref{eq:pg} at $\theta=\theta_{old}$, and it can be maximized with
several gradient steps on the same data. Left alone, though, it would keep
pushing $\rho_t$ for good actions towards infinity. PPO therefore
maximizes the \emph{clipped} objective
\begin{equation}
L^{\mathrm{clip}}(\theta)=\mathbb E_t\Big[\min\big(\rho_t\hat A_t,\;\mathrm{clip}(\rho_t,\,1-\epsilon,\,1+\epsilon)\,\hat A_t\big)\Big],\qquad \epsilon=0.2 .
\label{eq:clip}
\end{equation}
Two cases make it concrete. If $\hat A_t>0$ (a good action), the objective
grows with $\rho_t$ until $\rho_t=1.2$ and is flat beyond, so there is no
incentive to make that action more than 20\% more likely in one round. If
$\hat A_t<0$ (a bad action), the objective improves as $\rho_t$ falls, but
only until $\rho_t=0.8$. The $\min$ makes sure clipping can only remove
incentives, never add them. The result is a cheap, approximate trust
region.

The quantity actually minimized by gradient descent combines the actor and
critic losses:
\begin{equation}
L(\theta,\phi)=-L^{\mathrm{clip}}(\theta)+c_v\,\tfrac12\,\mathbb E_t\big[(V_\phi(o_t)-\hat R_t)^2\big],\qquad c_v=0.5 .
\label{eq:ppo_loss}
\end{equation}
The first term improves the policy and the second fits the critic by least
squares.

\subsection{The full training loop}
\begin{center}
\fbox{\begin{minipage}{0.92\textwidth}
\small
\textbf{PPO as implemented in \code{train\_ppo.py}}
\begin{enumerate}\itemsep1pt
\item Initialize actor $\theta$, critic $\phi$, $\log\sigma=-0.5$. Start 16 simulators, each with its own random season.
\item Repeat for 100 updates (single reservoir) or 245 updates (cascade):
\begin{enumerate}\itemsep0pt
\item \emph{Rollout:} each simulator runs 256 days with the current stochastic policy, so 4{,}096 $(o,a,r)$ records per update. A finished season immediately restarts with a new random season.
\item \emph{Advantages:} compute $\delta_t$, $\hat A_t$ (Eq.~\ref{eq:gae}) and $\hat R_t$ for every record; standardize $\hat A_t$.
\item \emph{Optimize:} 10 passes (epochs) over the 4{,}096 records in shuffled minibatches of 256; on each minibatch take one Adam gradient step on $L(\theta,\phi)$ (Eq.~\ref{eq:ppo_loss}), with gradient norm clipped at 0.5.
\end{enumerate}
\item Save $\theta$ (the trained policy).
\end{enumerate}
\end{minipage}}
\end{center}

\subsection{Hyperparameters and what they control}
\begin{center}
\small
\begin{tabular}{lrp{0.55\textwidth}}
\toprule
setting & value & what it controls\\
\midrule
parallel simulators & 16 & throughput; one OHQ instance per CPU worker\\
rollout length & 256 days & experience collected per simulator per update\\
samples per update & 4{,}096 & size of the batch the policy learns from each round\\
epochs per update & 10 & how many times each batch is reused\\
minibatch size & 256 & records per gradient step\\
learning rate (Adam) & $3\times10^{-4}$ & gradient step size\\
discount $\gamma$ & 0.99 & how far ahead the agent looks (effective horizon $\sim1/(1-\gamma)=100$ days)\\
GAE $\lambda$ & 0.95 & bias--variance trade-off of the advantage estimate\\
clip $\epsilon$ & 0.2 & maximum change in action probabilities per round (trust region)\\
value coefficient $c_v$ & 0.5 & weight of the critic's loss\\
gradient-norm clip & 0.5 & guards against occasional huge gradients\\
total steps & 409{,}600 / $\sim10^6$ & training budget, single reservoir / cascade\\
\bottomrule
\end{tabular}
\end{center}
These are standard PPO defaults. None were tuned for this problem.

\paragraph{Parallel rollouts.} The 16 simulators run in separate processes
(\path{gymnasium.vector.AsyncVectorEnv}). Each owns its own OHQ instance
and working folder. When a season ends, Gymnasium restarts it one step
later and returns one dummy record, which is masked out of the update.

\paragraph{A performance pitfall.} The first parallel run made almost no
progress. Every worker sat at $\approx370\%$ CPU because OHQ's solver
defaults to \code{n\_threads = 8} (OpenMP), and BLAS/PyTorch add their own
threads. Sixteen workers times eight threads oversubscribed the 32-core
machine. Whenever OHQ models run in parallel, \code{n\_threads} must be 1.
With \code{n\_threads = 1} in OHQ and \code{OMP\_NUM\_THREADS = 1} /
\code{torch.set\_num\_threads(1)} in the workers, throughput rose to
$\approx1{,}200$ environment steps per second.

